\documentclass[aps, pra, preprint]{revtex4-2}

\usepackage{mathtools, bm, amssymb, fixdif, derivative, physics3, microtype}
\usephxmodu{ab, ab.braket, operator}
\usepackage[mono = false]{libertine}
\usepackage[hidelinks, colorlinks]{hyperref}
\usepackage{zref-clever}
\makeatletter
\patchcmd{\frontmatter@abstract@produce}
  {\vskip200\p@\@plus1fil
   \penalty-200\relax
   \vskip-200\p@\@plus-1fil}{}{}{}
\let\oldHyPsd@CatcodeWarning\HyPsd@CatcodeWarning
\renewcommand \HyPsd@CatcodeWarning [1]%
  {\ifnum\pdfstrcmp{#1}{math shift}=0 \else \oldHyPsd@CatcodeWarning{#1} \fi}
\pdfstringdefDisableCommands{\let\HyPsd@CatcodeWarning\@gobble}
\renewcommand *\@email[2]{\endgroup \@AF@join
  {#1\href {mailto:#2}{\ttfamily #2}}}
\makeatother
\newtheorem{example}{Example}

\begin{document}

\title{MKCT: Note on Quantum Case v.s. Classical Case}
\author{Mingyu Xia}
\email{xiamingyu@westlake.edu.cn}
\affiliation{Department of Physics, Westlake University}
\date{Released 2026-mm-dd}

\maketitle

\section{Open quantum system}

The Hamiltonian for the model can be described by three terms,
\begin{equation}\label{eq:tot.hamiltonian}
  \hat H = \hat H_\text{sys} + \hat H_\text{bath} + \hat H_\text{sys-bath}
= \hat H_\text{sys}
  + \sum_j \frac{\hbar \omega_j}{2} (\hat p_j^2 + \hat q_j^2)
  + \hat V U(\hat q)
\end{equation}
where $\hat H_\text{sys}$ represents the system,
$\hat H_\text{bath}$ represents the bath, and
$\hat H_\text{sys-bath}$ represents their interaction.
The bath mode operators
\[
  \hat q_j \equiv \frac{\hat b_j^\dagger + \hat b_j}{\sqrt2}, \quad
  \hat p_j \equiv \frac{\hat b_j - \hat b_j^\dagger}{\iu\sqrt2}
\]
are the position and momentum operators for the $j$th mode of a harmonic bath
with frequency $\omega_j$,
which can be represented in ladder operators
\[
  \hat H_\text{bath}
= \sum_j \hbar \omega_j \ab(\hat b_j^\dagger \hat b_j + \frac12).
\]
The arbitrary potential energy expressed in polynomial
\[
  U(\hat q) = \alpha_0 + \alpha_1 \hat q + \alpha_2 \hat q^2 + \ldots,
\]
with $\hat q = \sum_j c_j \hat q_j$ stands for the collective mode.

The system term will be determined under different usages.
\emph{We take the God-given unit ($\hbar = 1$) for simplification in the
following derivations.}

\section{Application of the Liouvillian}

In the Application of the Liouvillian,
to compute $(\iu\mathcal L)^n\hat A$ to arbitrary order,
write down the first few terms, which satisfy
\begin{equation}\label{eq:act.Liouvillian}
  (\iu\mathcal L)^n\hat A = \sum_\text{all terms}
    \hat O \otimes \prod_j \hat Q_{m_j} \prod_k \hat P_{n_k},
\end{equation}
with the generalized bath mode operators
\begin{align*}
  \hat P_n = \sum_j c_j \omega_j^n \hat p_j,\quad
  \hat Q_m = \sum_j c_j \omega_j^m \hat q_j.
\end{align*}

Given the generator,
\begin{equation}\label{eq:generator}
  \hat G = \sum_\text{all terms} \hat O \otimes
    \prod_j \hat Q_{m_j} \prod_k \hat P_{n_k}
= \hat O \prod_j \hat Q_{m_j} \prod_k \hat P_{n_k},
\end{equation}
where a generic term in the sum is of the form
$\hat O \prod_j \hat Q_{m_j} \prod_k \hat P_{n_k}$;
each term is a product of a system operator $\hat O$
and a polynomial of bath operators.

\subsection{Quantum case}

In quantum case, acting the Liouvillian on the generator, we have
\[
  \iu\mathcal L\hat G = \iu[\hat H, \hat G],
\]
or more commonly,
\begin{equation}\label{eq:quantum.Liouvillian}
  \iu\mathcal L \mathord\cdot = \iu[\hat H, \mathord\cdot].
\end{equation}

The Liouvillian can be also split into three corresponding terms,
\begin{equation}\label{eq:quantum.Liouvillian.act}
    \iu\mathcal L \hat G =
    \iu\mathcal L_\text{sys}      \hat G
  + \iu\mathcal L_\text{bath}     \hat G
  + \iu\mathcal L_\text{sys-bath} \hat G
= \iu[\hat H_\text{sys},      \hat G]
+ \iu[\hat H_\text{bath},     \hat G]
+ \iu[\hat H_\text{sys-bath}, \hat G],
\end{equation}
which will be discussed respectively.

\subsubsection{The System term in quantum case}

Substituting $\hat G$ into the first (system) term,
it is straightforward that
\begin{equation}\label{eq:quantum.Liouvillian.sys}
  \iu\mathcal L_\text{sys} \hat G
= \iu[\hat H_\text{sys}, \hat G]
= \iu[\hat H_\text{sys}, \hat O] \prod_j \hat Q_{m_j} \prod_k \hat P_{n_k},
\end{equation}
which has also been demonstrated in~\cite{10.1063/5.0273707}.

\subsubsection{The bath term in quantum case}

For the second (bath) term, act the Liouvillian on $\hat p_j$ first
\begin{align*}
  \iu\mathcal L_\text{bath} \hat p_j & = \iu[\hat H_\text{bath}, \hat p_j]
= \frac{\iu}{2} \omega_j [\hat p_j^2 + \hat q_j^2, \hat p_j]
= -\omega_j \hat q_j,\\
\intertext{where the relation
  $[AB, C] = A[B, C] + [A, C]B$ is applied. Similarly, act it on $\hat q_j$}
  \iu\mathcal L_\text{bath} \hat q_j & = \iu[\hat H_\text{bath}, \hat q_j]
= \frac{\iu}{2} \omega_j [\hat p_j^2 + \hat q_j^2, \hat q_j]
= +\omega_j \hat p_j.
\end{align*}

Furthermore, act the Liouvillian on on the generalized coordinates
$\hat P_n$ and $\hat Q_m$, one can obtain
\begin{equation}\label{eq:quantum.Liouvillian.act.coordinate_momentum}
  \iu\mathcal L_\text{bath} \hat P_n
= \sum_j c_j \omega_j^n (\iu \mathcal L_\text{bath} \hat p_j)
= -\sum_j c_j \omega_j^{n+1} \hat q_j \equiv -\hat Q_{n+1},\quad
  \iu\mathcal L_\text{bath} \hat Q_m = +\hat P_{m+1}.
\end{equation}

Next, by substituting~\zcref{eq:quantum.Liouvillian.act.coordinate_momentum}
and $H_\text{bath}$ in~\zcref{eq:quantum.Liouvillian},
making use of the general commutation relation,
\[
  \ab[A, \prod_{i=1}^N B_i]
= \sum_{i=1}^N \prod_{\alpha_1=0}^{i-1} B_{\alpha_1} [A, B_i]
               \prod_{\alpha_2=i+1}^N B_{\alpha_2},\quad
  B_0 \equiv \identity,
\]
one can obtain the result of acting the Liouvillian on the bath term
in quantum case
\begin{equation}\label{eq:quantum.Liouvillian.act.bath}
  \begin{aligned}
  \iu\mathcal L_\text{bath} \hat G
& = \iu \hat O \ab[\hat H_\text{bath}, \prod_j \hat Q_{m_j}]
    \prod_k \hat P_{n_k}
    + \iu \hat O \prod_j \hat Q_{m_j}
    \ab[\hat H_\text{bath}, \prod_k \hat P_{n_k}]\\
& = \iu \hat O \sum_j
    \prod_{\alpha_1<j} \hat Q_{m_{\alpha_1}}
    [\hat H_\text{bath}, \hat Q_{m_j}]
    \prod_{\alpha_2>j} \hat Q_{m_{\alpha_2}}
    \prod_k \hat P_{n_k}\\
& \quad
    + \iu \hat O \prod_j \hat Q_{m_j} \sum_k
    \prod_{\alpha_1<k} \hat P_{n_{\alpha_1}}
    [\hat H_\text{bath}, \hat P_{n_k}]
    \prod_{\alpha_2>k} \hat P_{n_{\alpha_2}}\\
& = \hat O \Biggl(\sum_j
    \prod_{\alpha_1<j} \hat Q_{m_{\alpha_1}} \hat P_{m_j+1}
    \prod_{\alpha_2>j} \hat Q_{m_{\alpha_2}}
    \prod_k \hat P_{n_k}\\
& \quad - \prod_j \hat Q_{m_j} \sum_k
    \prod_{\alpha_1<k} \hat P_{n_{\alpha_1}} \hat Q_{n_k+1}
    \prod_{\alpha_2>k} \hat P_{n_{\alpha_2}}\Biggr),
  \end{aligned}
\end{equation}
which satisfies the expression in~\cite{10.1063/5.0273707}.
The derivation of acting the Liouvillian on the bath term
used $[\hat H_\text{bath}, \hat O] = 0$.

\subsubsection{The system-bath term in quantum case}

Since there exists
\[
  [\hat O, \hat q_j] = [\hat O, \hat p_j]
= [\hat V, \hat q_j] = [\hat V, \hat p_j] = 0, \quad [\hat O, \hat V] \neq 0,
\]
then acting the Liouvillian on the system-bath term directly,
\begin{equation}\label{eq:quantum.Liouvillian.act.sys-bath}
  \begin{aligned}
  \iu\mathcal L_\text{sys-bath} \hat G = \iu[\hat H_\text{sys-bath}, \hat G]
& = \iu\hat V U(\hat q) \hat O \prod_j \hat Q_{m_j} \prod_k \hat P_{n_k} -
    \iu\hat O \prod_j \hat Q_{m_j} \prod_k \hat P_{n_k} \hat V U(\hat q)\\
& = \iu \hat V \hat O U(\hat q) \prod_j \hat Q_{m_j} \prod_k \hat P_{n_k}
  - \iu \hat O \hat V \prod_j \hat Q_{m_j} \prod_k \hat P_{n_k} U(\hat q)
  \end{aligned}
\end{equation}
which satisfies the expression in~\cite{10.1063/5.0273707}.

\subsection{Classical case}

In the classical case, the bath modes are described by canonical coordinates
$q_j$ and momenta $p_j$ with Poisson bracket $\{q_j, p_k\} = \delta_{jk}$,
all products are ordinary commutative multiplication.

The system Hamiltonian depends only on the system variables,
the bath and system-bath Hamiltonians are
\[
  H_\text{bath} = \sum_j \frac{\omega_j}{2}(p_j^2 + q_j^2),\quad
  H_\text{sys-bath} = V U(q),\quad q = \sum_j c_j q_j,
\]
where $V$ is a function of the system variables and $U(q)$ is a smooth function
of the collective bath coordinate.

The generator (a generic term) is now a classical function
\[
  G = O \prod_j Q_{m_j} \prod_k P_{n_k},\quad
  Q_m = \sum_j c_j \omega_j^m q_j,\quad
  P_n = \sum_j c_j \omega_j^n p_j,
\]
with $O$ a function of the system variables.

In contrast to the quantum case,
the Liouvillian in \zcref{eq:quantum.Liouvillian} will become
\begin{equation}
  \iu\mathcal L \cdot = \{\cdot, H\}
= \sum_i \ab(\pdv{\,\mathord\cdot}{q_i} \pdv H {p_i}
  - \pdv H {q_i} \pdv{\,\mathord\cdot}{p_i}),
\end{equation}
similarly, it can be splited into three corresponding terms,
\begin{equation}\label{eq:classical.Liouvillian.act}
  \iu\mathcal LG = \iu\mathcal L_\text{sys}G
  + \iu\mathcal L_\text{bath}G + \iu\mathcal L_\text{sys-bath}G.
\end{equation}
which will be discussed respectively.

\subsubsection{The system term in classical case}

Substituting $\hat G$ into the first (system) term,
it is straightforward that
\begin{equation}\label{eq:classical.Liouvillian.sys}
  \iu\mathcal L_\text{sys}G
  = \{O, H_\text{sys}\} \prod_j Q_{m_j} \prod_k P_{n_k}.
\end{equation}

\subsubsection{The bath term in classical case}

For the second (bath) term, act the Liouvillian on $q_j$ and $p_j$
respectively first
\[
  \iu \mathcal L_\text{bath} q_j = \{q_j, H_\text{bath}\} = \omega_j p_j,\quad
  \iu \mathcal L_\text{bath} p_j = \{p_j, H_\text{bath}\} = -\omega_j q_j,
\]
where the filtering property of the $\delta$-function
$\sum_i \pdv{x_j}{x_i} f_i = \sum_i \delta_{ji} f_i = f_j$ is applied.

Furthermore, act the Liouvillian on the generalized coordinates
$P_n$ and $Q_m$, one can obtain
\begin{equation}\label{eq:classical.Liouvillian.act.coordinate_momentum}
  \iu\mathcal L_\text{bath} Q_m
  = \{Q_m, H_\text{bath}\} = P_{m+1},\quad
  \{P_n, H_\text{bath}\} = -Q_{n+1}.
\end{equation}

Since the Poisson brackets satisfy an identity,
$\{fh, g\} = \{f, g\}h + f\{h, g\}$, which is similar to the commutator,
then, we can also obtain the general product rule,
\[
  \ab\{\prod_{i=1}^N g_n, f\}
= \sum_{i=1}^N \{g_i, f\} \prod_{\alpha=0,\,\alpha\neq i}^N g_\alpha,\quad
  g_0 = 1,
\]
one can obtain the result of acting the Liouvillian on the bath term
in classical case
\begin{equation}\label{eq:classical.Liouvillian.act.bath}
  \begin{aligned}
  \iu\mathcal L_\text{bath} G
& = O \ab\{\prod_j Q_{m_j}, H_\text{bath}\} \prod_k P_{n_k}
  + O \prod_j Q_{m_j} \ab\{\prod_k P_{n_k}, H_\text{bath}\}\\
& = O \sum_j \{Q_{m_j}, H_\text{bath}\}
    \prod_{\alpha\neq j} Q_{m_\alpha} \prod_k P_{n_k}
    + O \prod_j Q_{m_j} \sum_k \{P_{n_k}, H_\text{bath}\}
    \prod_{\alpha\neq k} P_{n_\alpha}\\
& = O \ab[\sum_j P_{m_{j+1}} \prod_{\alpha\neq j} Q_{m_\alpha} \prod_k P_{n_k}
        - \sum_k Q_{n_{k+1}} \prod_j Q_{m_j} \prod_{\alpha\neq k} P_{n_\alpha}].
  \end{aligned}
\end{equation}
The result is much more elegant than that of the quantum case,
since the order of $P$ and $Q$ can be exchanged freely.

\subsubsection{The system-bath term in classical case}

Since only the Poisson brackets of $q$, $O$, $V$ to $U(q)$ are zero,
we only have to expanding the Poisson bracket
of the generalized bath mode function using the general product rule,
\begin{align*}
  \ab\{P_{n_k}, U(q)\}
& = \ab\{\sum_j c_j \omega_j^{n_k} p_j,
      \sum_{i=0} \alpha_i \sum_{j'} c_{j'}^i q_{j'}^i\}\\
& = -\sum_j c_j \omega_j^{n_k} \sum_{i=1} \alpha_i \sum_{j'}
      c_{j'}^i \sum_{i'} i q_{j'}^{i-1} \delta_{j'i'} \delta_{ji'}\\
& =-\sum_j c_j^2 \omega_j^{n_k} \sum_{i=1} i \alpha_i c_j^{i-1} q_j^{i-1}
= -\odv{U(q)}{q} \sum_j c_j^2 \omega_j^{n_k} = -U'(q) \mu_{n_k},
\end{align*}
where $\mu_n \equiv \sum_j c_j^2 \omega_j^n$ denotes the bath spectral moments.
Then, using the general product rule,
\[
  \ab\{\prod_k P_{n_k}, U(q)\}
= -U'(q) \sum_k \mu_{n_k} \prod_{\alpha\neq k} P_{n_\alpha}.
\]

Finally, acting the Liouvillian on the third (system-bath) term
by substituting the result above,
\begin{equation}\label{eq:classical.Liouvillian.act.sys-bath}
  \begin{aligned}
    \iu\mathcal L_\text{sys-bath}G & = \{G, H_\text{sys-bath}\}\\
& = \{O, V\} U(q) \prod_j Q_{m_j} \prod_k P_{n_k}
    + O V \prod_j Q_{m_j} \ab\{\prod_k P_{n_k}, U(q)\},\\
& = \{O, V\} U(q) \prod_j Q_{m_j} \prod_k P_{n_k}
    - O V U'(q) \prod_j Q_{m_j}
    \sum_k \mu_{n_k} \prod_{\alpha\neq k} P_{n_\alpha}.
  \end{aligned}
\end{equation}
Then, we can evaluate the higher-order moments $\Omega_n$ in
Memory kernel coupling theory (MKCT).

\section{Evaluation of the Mori products}

In~\cite{10.1063/5.0273707}, the Mori-product has been defined as
\begin{equation}\label{eq:quantum.mori.product}
  (\hat A, \hat B) =
  \Tr(\hat A\hat B^\dagger \hat\sigma_0 \otimes \hat\rho_\text{eq}^\text{bath})
\end{equation}
for two arbitrary system operators $\hat A$ and $\hat B$,
$\hat\sigma_0$ is the initial system density operator,
and $\hat\rho_\text{eq}^\text{bath}$ is the equilibrium bath density operator,
with $\beta = 1/T$ being the inverse temperature.

As the instruction in~\cite{10.1063/5.0273707},
we need to evaluate the Mori-product of the generator $\hat G$
and an arbitrary system operator $\hat A$, it can be expressed as
\begin{equation}\label{eq:quantum.generator.mori-product}
  (\hat G, \hat A)
= \Tr(\hat O \hat A \hat\sigma_0)
  \Tr\ab(\prod_j \hat Q_{m_j} \prod_k \hat P_{n_k}
         \hat \rho_\text{eq}^\text{bath}),
\end{equation}
this evaluation is in the quantum case.
Now, we evaluate the product in the classical case.

In the classical case, the Mori product for two phase-space functions $X$
and $Y$ is defined as the average over the initial factorised distribution:
\[
  (X, Y) = \int \d\Gamma \rho_0(\Gamma) X(\Gamma) Y(\Gamma),\quad
  \rho_0 = \sigma_0(\text{sys}) \times \rho_\text{eq}^\text{bath}(\text{bath}).
\]

Here $\sigma_0$ is an arbitrary initial system distribution
(e.g., a pure state or a thermal state) and $\rho_\text{eq}^\text{bath}$
is the classical canonical distribution of the harmonic bath:
\[
  \rho_\text{eq}^\text{bath}
  = \prod_j \frac{\beta\omega_j}{2\pi} \upe^{-\beta\omega_j(p_j^2+q_j^2)/2},
\]
with $\beta = 1/T$.
The bath modes are independent Gaussian variables with zero mean
and variances $\expval*{q_j^2} = \expval*{p_j^2} = 1/(\beta\omega_j)$.

\subsection{Factorisation of the Mori product}

Consider a generic term of the generator
(as obtained from repeated action of the Liouvillian)
\[
  G = O(\text{sys}) \prod_j Q_{m_j} \prod_k P_{n_k},
\]
where $O$ depends only on the system variables, and
\[
  Q_m = \sum_j c_j \omega_j^m q_j,\quad
  P_n = \sum_j c_j \omega_j^n p_j.
\]

For an arbitrary system function $A$ (depending only on system variables),
the Mori product factorises because the system and bath distributions
are independent:
\[
  (G, A) = \ab(\int \d\text{sys} \sigma_0(\text{sys})
               O(\text{sys}) A(\text{sys})) \times
  \expval{\prod_j Q_{m_j} \prod_k P_{n_k}}_\text{bath},
\]
where $\expval{\cdot}_\text{bath}$ denotes
the average over $\rho_\text{eq}^\text{bath}$.

\subsection{Bath averages}

Because the bath equilibrium is a product of independent Gaussians
for each mode, the variables $\{q_j\}$ and $\{p_j\}$ are independent.
Consequently, the sets $\{Q_m\}$ and $\{P_n\}$ are independent as well,
and each set is jointly Gaussian with zero mean. Their covariance matrices are
\[
  \expval{Q_m Q_n} = \sum_j c_j^2 \omega_j^{m+n} \expval{q_j^2}
= \frac{1}{\beta} \sum_j c_j^2 \omega_j^{m+n-1},\quad
  \expval{P_m P_n} = \frac{1}{\beta} \sum_j c_j^2 \omega_j^{m+n-1}
= \expval{Q_m Q_n},
\]
and $\expval{Q_m P_n} = 0$.

Introduce the classical bath spectral moments
\[
  \mu_n = \sum_j c_j^2 \omega_j^n,\quad
  \eta_n^{\text{class}} = \frac{1}{\beta} \mu_{n-1}
= \frac{2}{\pi\beta}\int_0^\infty \d\omega J(\omega) \omega^{n-1},
\]
where $J(\omega) = \frac{\pi}{2}\sum_j c_j^2 \delta(\omega-\omega_j)$
is the spectral density. Then
\[
  \expval{Q_m Q_n} = \eta_{m+n}^{\text{class}},\quad
  \expval{P_m P_n} = \eta_{m+n}^{\text{class}}.
\]

\subsection{Generating function}

To compute the average of a product of several $Q$'s and $P$'s,
we use the Gaussian generating function. For the $Q$'s,
\[
\expval{\exp\ab(\sum_a \lambda_a Q_{m_a})}_\text{bath}
= \exp\ab(\frac12 \sum_{a,b} \lambda_a\lambda_b \expval{Q_{m_a}Q_{m_b}}),
\]
and similarly for the $P$'s with independent variables $\{\lambda'_c\}$.
Because $Q$ and $P$ are independent, the joint generating function factorises:
\[
  f(\{\lambda_a\},\{\lambda'_c\}) =
  \exp\ab(\frac12 \sum_{a,b} \lambda_a\lambda_b \eta_{m_a+m_b}^{\text{class}})
  \exp\ab(\frac12 \sum_{c,d} \lambda'_c\lambda'_d
    \eta_{n_c+n_d}^{\text{class}}).
\]

The desired average is then
\[
  \expval{\prod_j Q_{m_j} \prod_k P_{n_k}}_\text{bath}
= \ab(\prod_j \frac{\partial}{\partial\lambda_j})
  \ab(\prod_k \frac{\partial}{\partial\lambda'_k})
  f(\{\lambda\},\{\lambda'\}) \Big|_{\lambda=\lambda'=0}.
\]
Because the generating function is even, the average vanishes
unless the number of $Q$ factors and the number of $P$ factors are both even.
When they are even, the result factorises into a product
of two independent Gaussian moments:
\[
  \expval{\prod_j Q_{m_j}}_\text{bath}
= \sum_\text{pairings of $\{j\}$}
  \prod_\text{pairs $(a,b)$} \eta_{m_a+m_b}^{\text{class}},\quad
  \expval{\prod_k P_{n_k}}_\text{bath}
= \sum_{\text{pairings of }\{k\}}
  \prod_{\text{pairs }(c,d)} \eta_{n_c+n_d}^{\text{class}},
\]
and the full average is the product of these two.

\subsection{System part}

The system part is simply the correlation of $O$ and $A$
under the initial system distribution:
\[
  \int \d\text{sys} \sigma_0(\text{sys}) O(\text{sys}) A(\text{sys}).
\]

This can be evaluated analytically or numerically depending on the system.

Finally, we obtain the evaluation in the classical case
\begin{equation}\label{eq:classical.generator.mori-product}
  (\hat G, \hat A)
= \ab(\int \d \text{sys} \sigma_0(\text{sys}) OA) \times
  \ab(\expval{\prod_j Q_{m_j}}_\text{bath}) \times
  \ab(\expval{\prod_j P_{n_k}}_\text{bath})
\end{equation}
with the bath averages given by the pair-sum formulas above.
This is the classical analogue of Eq. (29) in the quantum paper.
The absence of cross terms between $Q$ and $P$
and the simplicity of Gaussian moments
make the classical evaluation significantly more straightforward.

\bibliography{reference}

\end{document}
